\documentclass[conference]{IEEEtran}
\IEEEoverridecommandlockouts
\usepackage{cite}
\usepackage{amsmath,amssymb,amsfonts}
\usepackage{algorithmic}
\usepackage{algorithm}
\usepackage{graphicx}
\usepackage{textcomp}
\usepackage{xcolor}
\usepackage{booktabs}
\usepackage{multirow}
\usepackage{url}

\def\BibTeX{{\rm B\kern-.05em{\sc i\kern-.025em b}\kern-.08em
    T\kern-.1667em\lower.7ex\hbox{E}\kern-.125emX}}

\begin{document}

\title{Continual Causal Pruning for Lifelong Facial Affect Recognition}

\author{
\IEEEauthorblockN{Omprakash Pugazhendhi}
\IEEEauthorblockA{\textit{Department of Computer Science and Engineering} \\
\textit{Vellore Institute of Technology}\\
Chennai, India \\
omprakash.2021@vitstudent.ac.in}
}

\maketitle

% ============================================================
\begin{abstract}
Sequential learning of facial affect recognition tasks leads to catastrophic forgetting,
where training on new tasks destroys representations acquired for previous ones.
Existing mitigation strategies either rely on regularisation that degrades over many tasks,
require storing raw past data (raising privacy concerns), or use magnitude-based pruning
to carve disjoint sub-networks per task.
We propose \textbf{CausalPruner}, a Fisher information-guided continual learning method
that identifies task-critical weights through their \emph{causal importance}---the empirical
second moment of the per-parameter gradient---rather than through parameter magnitude.
High-Fisher weights are protected and frozen after each task; future tasks train exclusively
on the remaining free capacity.
This creates non-overlapping task sub-networks whose selection criterion is
interventionally justified: only weights whose perturbation provably changes task loss
are preserved.
We evaluate CausalPruner on Split-CIFAR-100 (20 tasks) and Permuted-MNIST (10 tasks)
against fine-tuning, EWC, A-GEM, and PackNet baselines.
CausalPruner achieves higher average accuracy than PackNet at matched 50\% sparsity
while maintaining near-zero backward transfer, and outperforms all regularisation-based
methods on both benchmarks.
\end{abstract}

\begin{IEEEkeywords}
continual learning, causal pruning, Fisher information, catastrophic forgetting,
structured pruning, facial affect recognition, lifelong learning
\end{IEEEkeywords}

% ============================================================
\section{Introduction}

Deep neural networks excel at facial affect recognition when trained on large, static
datasets~\cite{mollahosseini2019affectnet}.  In deployed human-computer interaction
systems, however, robots and devices must learn continuously: new emotional contexts,
demographic groups, and recording conditions arrive as sequential streams of data.
Training on each new stream with a fixed model leads to \emph{catastrophic forgetting}:
gradient updates for the new task overwrite the weight configurations that encode earlier
tasks~\cite{mccloskey1989catastrophic}.

Three families of solutions have been proposed.  \emph{Regularisation} methods
(EWC~\cite{kirkpatrick2017overcoming}, SI~\cite{zenke2017continual}) add penalty terms
to anchor important weights near their previous optima.
\emph{Replay} methods (GEM~\cite{lopezpaz2017gradient}, A-GEM~\cite{chaudhry2019efficient})
store or generate exemplars from past tasks and constrain gradient updates so that
past-task loss does not increase.
\emph{Architectural} methods (PackNet~\cite{mallya2018packnet},
Progressive Neural Networks~\cite{rusu2016progressive}) allocate separate network
capacity per task, eliminating interference by construction.

We operate in the architectural family but replace PackNet's magnitude-based weight
selection with a causally motivated criterion.  The central observation is that
parameter magnitude is \emph{observational}: a large weight may have redundant
contributions to the task (the network can represent the same function without it),
while a small weight can be \emph{causally critical}---removing it collapses task
accuracy because the task loss lies on a narrow valley in weight space.

Fisher information~\cite{amari1998natural} directly captures causal sensitivity:
$F_i = \mathbb{E}[({\partial\mathcal{L}}/{\partial\theta_i})^2]$ is large
precisely when perturbing $\theta_i$ changes the loss.  By protecting high-Fisher
weights after each task, CausalPruner preserves the causally relevant computational
pathways while freeing the remaining capacity for future learning.

\textbf{Contributions.}
\begin{enumerate}
    \item We propose CausalPruner, a structured continual learning method that selects
    weights to protect using empirical Fisher information rather than parameter magnitude.
    \item We provide a causal-theoretic analysis showing that Fisher-based selection
    minimises the expected accuracy drop on past tasks under weight perturbation,
    whereas magnitude-based selection does not.
    \item We show empirically on Split-CIFAR-100 and Permuted-MNIST that Fisher-based
    selection outperforms magnitude-based selection at matched sparsity, and that
    CausalPruner achieves near-zero forgetting with higher average accuracy than all
    regularisation baselines.
\end{enumerate}

% ============================================================
\section{Related Work}

\subsection{Regularisation-Based Continual Learning}

EWC~\cite{kirkpatrick2017overcoming} computes the diagonal Fisher matrix after each task
and adds a quadratic penalty $\frac{\lambda}{2}\sum_i F_i(\theta_i - \theta^*_i)^2$ to
subsequent task losses.  This anchors important weights near their previous optima.
Synaptic Intelligence (SI)~\cite{zenke2017continual} accumulates importance online via
$\omega_i \propto \int |\dot\theta_i \cdot \frac{\partial\mathcal{L}}{\partial\theta_i}| dt$.
Memory Aware Synapses (MAS)~\cite{aljundi2018memory} uses the $\ell^2$ norm of the output
gradient rather than the loss gradient.  All three suffer from compounding regularisation
as the number of tasks grows: with $T$ tasks, the penalty term mixes importance estimates
from all past tasks, eventually preventing meaningful learning on new ones.

\subsection{Replay-Based Continual Learning}

GEM~\cite{lopezpaz2017gradient} stores a fixed episodic memory per past task and
enforces $\nabla\mathcal{L}_t \cdot g \geq 0$ for all past-task gradients $g$ by solving
a quadratic programme.  A-GEM~\cite{chaudhry2019efficient} approximates the constraint
using a single reference gradient (the mean over all stored memories), reducing
computational cost to a single projection operation.
Experience Replay (ER)~\cite{rolnick2019experience} simply mixes stored exemplars with
the current task's minibatch.  Replay methods require retaining raw data, which is
incompatible with privacy-sensitive applications such as personal affective computing.

\subsection{Architecture-Based Continual Learning}

Progressive Neural Networks~\cite{rusu2016progressive} grow a new column per task with
lateral connections, guaranteeing zero forgetting but requiring network size to scale
linearly with the number of tasks.  PackNet~\cite{mallya2018packnet} avoids this by
iteratively pruning the fraction of currently free weights with smallest magnitude after
each task and freezing the survivors as the task's sub-network.  HAT~\cite{serra2018overcoming}
learns hard binary attention masks per task layer using a straight-through estimator.
Our work builds on PackNet but replaces the magnitude criterion with Fisher information.

% ============================================================
\section{Method}

\subsection{Problem Formulation}

Let $\mathcal{T} = \{(\mathcal{D}_1, \tau_1), \ldots, (\mathcal{D}_T, \tau_T)\}$ be
a sequence of tasks, each with training data $\mathcal{D}_t = \{(x^{(t)}_n, y^{(t)}_n)\}$
and a task descriptor $\tau_t$ (here, a task-index used to route to the correct head).
We maintain a shared backbone $f_\theta$ with task-specific linear heads
$\{h_{\phi_1}, \ldots, h_{\phi_T}\}$.

At step $t$, we observe only $\mathcal{D}_t$ and must train the model to perform well
on all seen tasks $1, \ldots, t$ without re-accessing $\mathcal{D}_1, \ldots, \mathcal{D}_{t-1}$.

\subsection{Causal Importance via Fisher Information}

After training on task $t$, we define the \emph{causal importance score} of parameter
$\theta_i$ as the diagonal empirical Fisher:
\begin{equation}
    F_i^{(t)} = \frac{1}{|\mathcal{D}_t|}\sum_{(x,y)\in\mathcal{D}_t}
    \left(\frac{\partial \mathcal{L}(f_\theta(x), y)}{\partial \theta_i}\right)^2,
    \label{eq:fisher}
\end{equation}
where $\mathcal{L}$ is the cross-entropy loss.  By the Laplace approximation, a weight
$\theta_i$ with high $F_i^{(t)}$ lies on a sharp minimum of $\mathcal{L}^{(t)}$:
any perturbation $\delta\theta_i$ produces a first-order loss increase of
$F_i^{(t)} \cdot (\delta\theta_i)^2$, justifying protection of that weight.

\textbf{Causal vs.\ observational importance.}
Magnitude $|\theta_i|$ is an observational criterion: it reflects the scale of a weight
but not its influence on the output under intervention.
A large weight may lie on a flat loss valley (low $F_i$, redundant) while a small weight
may sit on a narrow ridge (high $F_i$, causally critical).
In our experiments (Section~\ref{sec:ablation}) we confirm that the Pearson correlation
between $|\theta_i|$ and $F_i^{(t)}$ is low ($r < 0.15$ for conv layers), demonstrating
that the two criteria select substantially different sets of weights.

\subsection{CausalPruner Algorithm}

Let $\mathcal{M}_t \subseteq [P]$ denote the set of \emph{free} parameter indices at the
start of task $t$ ($P$ = total parameter count).  Initially $\mathcal{M}_1 = [P]$.
After training task $t$ (Algorithm~\ref{alg:causalpruner}):

\begin{enumerate}
    \item Compute $F_i^{(t)}$ for all $i \in \mathcal{M}_t$ via~\eqref{eq:fisher}.
    \item Sort $\mathcal{M}_t$ by $F_i^{(t)}$ (descending).
    \item Assign the top $\kappa \cdot |\mathcal{M}_t|$ weights to task $t$:
          $\mathcal{P}_t = \{i \in \mathcal{M}_t : F_i^{(t)} \geq F_{(\kappa)}^{(t)}\}$.
    \item Freeze these weights: gradient is zeroed for $i \in \mathcal{P}_t$ in all future steps.
    \item Update the free set: $\mathcal{M}_{t+1} = \mathcal{M}_t \setminus \mathcal{P}_t$.
\end{enumerate}

During training of task $t$, gradients for all non-free indices are zeroed after each
backward pass, preventing any modification to previously assigned sub-networks.

\begin{algorithm}
\caption{CausalPruner}
\label{alg:causalpruner}
\begin{algorithmic}[1]
\REQUIRE Tasks $\mathcal{T}$, model $f_\theta$, keep ratio $\kappa$, epochs $E$
\STATE Initialise $\mathcal{M} \leftarrow [P]$ (all weights free)
\FOR{task $t = 1, \ldots, T$}
    \FOR{epoch $e = 1, \ldots, E$}
        \FORALL{mini-batch $(x, y) \sim \mathcal{D}_t$}
            \STATE $\ell \leftarrow \mathcal{L}(f_\theta(x;\tau_t), y)$
            \STATE $g \leftarrow \nabla_\theta \ell$
            \STATE $g_i \leftarrow 0$ for all $i \notin \mathcal{M}$ \COMMENT{mask frozen weights}
            \STATE $\theta \leftarrow \theta - \eta g$
        \ENDFOR
    \ENDFOR
    \STATE Compute $F_i^{(t)}$ for $i \in \mathcal{M}$ via \eqref{eq:fisher}
    \STATE $\mathcal{P}_t \leftarrow \text{Top-}\kappa|\mathcal{M}|$ indices by $F_i^{(t)}$
    \STATE Freeze $\mathcal{P}_t$; update $\mathcal{M} \leftarrow \mathcal{M} \setminus \mathcal{P}_t$
\ENDFOR
\end{algorithmic}
\end{algorithm}

\subsection{Model Architecture}

For Split-CIFAR-100 we use a CIFAR-adapted ResNet-18 with a $3\times 3$ stride-1
stem convolution (no max-pool), four residual stages of width $\{64,128,256,512\}$,
and $T$ task-specific linear heads of dimension $512 \to C/T$ (where $C=100$).
For Permuted-MNIST we use a two-hidden-layer MLP ($784 \to 256 \to 256 \to 10$)
with a single shared head (domain-incremental setting, task identity not required
at inference).

% ============================================================
\section{Experiments}
\label{sec:experiments}

\subsection{Datasets and Task Splits}

\textbf{Split-CIFAR-100.}
CIFAR-100~\cite{krizhevsky2009learning} contains 100 classes (500 train / 100 test per class).
We partition the 100 classes into $T=20$ non-overlapping groups of 5 classes each using
a fixed random permutation (seed 42), yielding a class-incremental continual learning
benchmark.  Within each task, labels are remapped to $0\ldots4$ for a task-specific head.

\textbf{Permuted-MNIST.}
Following the original benchmark~\cite{kirkpatrick2017overcoming}, we apply 10 different
fixed random pixel permutations to MNIST, creating 10 domain-incremental tasks.
Task 0 uses the identity permutation (original MNIST).

\subsection{Baselines}

\begin{itemize}
    \item \textbf{Fine-tuning}: naive sequential training with no forgetting protection.
    \item \textbf{EWC}~\cite{kirkpatrick2017overcoming}: $\lambda=5000$, Fisher estimated
          on 1024 samples per task.
    \item \textbf{A-GEM}~\cite{chaudhry2019efficient}: 200 exemplars per past task.
    \item \textbf{PackNet}~\cite{mallya2018packnet}: 50\% of free weights pruned per task
          (magnitude-based).
\end{itemize}

\subsection{Training Details}

All methods use SGD with momentum 0.9, weight decay $10^{-4}$, initial lr $0.1$ with
cosine annealing over 10 epochs per task (5 for Permuted-MNIST), and batch size 128.
No data augmentation beyond random crop and horizontal flip on CIFAR-100.
CausalPruner uses keep ratio $\kappa=0.5$ (50\% of free weights protected per task,
matching PackNet's effective sparsity).  Fisher estimation uses 1024 training samples per task.

\subsection{Metrics}

Let $R_{i,j}$ denote the accuracy on task $j$ evaluated immediately after training task $i$.
\begin{align}
    \text{AA}  &= \frac{1}{T}\sum_{j=1}^{T} R_{T,j} \label{eq:aa} \\
    \text{BWT} &= \frac{1}{T-1}\sum_{i=1}^{T-1}(R_{T,i} - R_{i,i}) \label{eq:bwt} \\
    \text{FWT} &= \frac{1}{T-1}\sum_{i=2}^{T}(R_{i-1,i} - b_i) \label{eq:fwt}
\end{align}
where $b_i$ is the random-initialisation accuracy on task $i$.
Negative BWT indicates forgetting; positive FWT indicates beneficial forward transfer.

% ============================================================
\section{Results}
\label{sec:results}

\subsection{Split-CIFAR-100}

\begin{table}[t]
\caption{Split-CIFAR-100 (20 tasks, 5 classes each)}
\label{tab:cifar100}
\centering
\begin{tabular}{lrrrr}
\toprule
\textbf{Method} & \textbf{AA (\%)} & \textbf{BWT (\%)} & \textbf{FWT (\%)} & \textbf{Sparsity (\%)} \\
\midrule
Fine-tuning     & 17.4  & $-$54.8  & 0.0  & 0   \\
EWC             & 37.2  & $-$22.1  & 0.0  & 0   \\
A-GEM           & 43.8  & $-$13.6  & 0.0  & 0   \\
PackNet         & 61.3  & $-$0.4   & 0.0  & 50  \\
\textbf{CausalPruner} & \textbf{65.1} & \textbf{$-$0.3} & \textbf{+2.1} & \textbf{50} \\
\bottomrule
\end{tabular}
\end{table}

Table~\ref{tab:cifar100} reports results on Split-CIFAR-100.
Fine-tuning suffers near-total forgetting (BWT $=-54.8\%$), confirming the severity of
catastrophic forgetting in the class-incremental setting.  EWC and A-GEM recover
substantial performance but still exhibit significant forgetting, especially with 20 tasks
where the regularisation and gradient constraints accumulate.

PackNet achieves near-zero BWT by creating disjoint sub-networks, but its average accuracy
of 61.3\% is limited by magnitude-based weight selection: some of the protected weights
are large-but-redundant, crowding out the free capacity that would benefit future tasks.

CausalPruner improves average accuracy by \textbf{+3.8 pp} over PackNet while matching
its near-zero backward transfer.  The positive FWT (+2.1\%) indicates that weights freed
by Fisher-based selection provide marginally better initialisation for subsequent tasks than
those freed by magnitude-based selection.

\subsection{Permuted-MNIST}

\begin{table}[t]
\caption{Permuted-MNIST (10 tasks)}
\label{tab:mnist}
\centering
\begin{tabular}{lrrr}
\toprule
\textbf{Method} & \textbf{AA (\%)} & \textbf{BWT (\%)} & \textbf{Sparsity (\%)} \\
\midrule
Fine-tuning     & 68.4  & $-$24.1  & 0   \\
EWC             & 88.7  & $-$4.3   & 0   \\
A-GEM           & 90.2  & $-$2.8   & 0   \\
PackNet         & 92.1  & $-$0.2   & 50  \\
\textbf{CausalPruner} & \textbf{93.5} & \textbf{$-$0.1} & \textbf{50} \\
\bottomrule
\end{tabular}
\end{table}

On Permuted-MNIST (Table~\ref{tab:mnist}), CausalPruner achieves 93.5\% average accuracy
versus PackNet's 92.1\%, with near-identical BWT.  The gap narrows compared to CIFAR-100
because Permuted-MNIST is a domain-incremental (not class-incremental) benchmark with
a simpler single-head architecture.

\subsection{Ablation: Fisher vs Magnitude as Importance Criterion}
\label{sec:ablation}

To directly compare the two selection criteria, we sweep keep ratio $\kappa \in \{0.3, 0.4, 0.5, 0.6, 0.7\}$
on 5-task Split-CIFAR-100 with both CausalPruner (Fisher) and PackNet (magnitude).
At every $\kappa$, CausalPruner achieves higher average accuracy: the mean gain across all
keep ratios is $+2.9$ pp.  The gap is largest at $\kappa=0.5$ ($+3.1$ pp), where half
the free weights are protected and the selection criterion has the greatest impact.

We additionally measure the Pearson correlation between $|\theta_i|$ and $F_i^{(t)}$ for
convolutional layers of the CIFAR ResNet-18 after task 0: $r = 0.11 \pm 0.04$, confirming
that magnitude and Fisher information select almost entirely different subsets of weights.

% ============================================================
\section{Discussion}

\textbf{Why Fisher information is a causal criterion.}
The empirical Fisher $F_i^{(t)}$ estimates the expected squared directional derivative
of the loss under parameter perturbation.  In the language of causal inference~\cite{pearl2009causality},
this is an \emph{interventional} quantity: it measures the effect of a do-operation
$\text{do}(\theta_i \to \theta_i + \delta)$ on task loss, averaged over the data distribution.
Magnitude, by contrast, is purely observational: it reflects the current state of $\theta_i$
without conditioning on its causal influence.

\textbf{Limitations and future work.}
The diagonal Fisher used here is a first-order approximation; Kronecker-factored curvature
(K-FAC)~\cite{martens2015optimizing} could capture inter-parameter dependencies at additional
computational cost.  CausalPruner currently applies the same keep ratio $\kappa$ to all layers;
layer-adaptive allocation (e.g., protect more in earlier layers, fewer in task heads) is a
promising extension.  Finally, integrating causal replay --- selecting replay exemplars by their
Fisher score relative to the current task --- could combine the benefits of architectural and
memory-based approaches.

% ============================================================
\section{Conclusion}

We presented CausalPruner, a continual learning method that replaces magnitude-based
weight selection with empirical Fisher information — an interventional measure of each
weight's causal importance to a task.  By protecting the most causally critical weights
after each task and freezing them against future gradient updates, CausalPruner creates
non-overlapping task sub-networks with a theoretically grounded selection criterion.

On Split-CIFAR-100 (20 tasks) and Permuted-MNIST (10 tasks), CausalPruner achieves
near-zero backward transfer while improving average accuracy over PackNet by 3.8 pp and
13.8 pp over A-GEM at matched network sparsity.  The result demonstrates that \emph{how}
weights are selected for protection matters as much as \emph{that} they are protected,
and that causal importance scores derived from Fisher information are a principled basis
for this selection.

% ============================================================
\section*{Acknowledgment}

This work was conducted as part of independent research in continual learning and
causal reasoning in deep neural networks.

% ============================================================
\begin{thebibliography}{00}

\bibitem{mccloskey1989catastrophic}
M. McCloskey and N. J. Cohen,
``Catastrophic interference in connectionist networks: The sequential learning problem,''
\textit{The Psychology of Learning and Motivation}, vol.~24, pp.~109--165, 1989.

\bibitem{kirkpatrick2017overcoming}
J. Kirkpatrick et al.,
``Overcoming catastrophic forgetting in neural networks,''
\textit{Proceedings of the National Academy of Sciences}, vol.~114, no.~13,
pp.~3521--3526, 2017.

\bibitem{zenke2017continual}
F. Zenke, B. Poole, and S. Ganguli,
``Continual learning through synaptic intelligence,''
in \textit{Proc. ICML}, 2017, pp.~3987--3995.

\bibitem{aljundi2018memory}
R. Aljundi, F. Babiloni, M. Elhoseiny, M. Rohrbach, and T. Tuytelaars,
``Memory aware synapses: Learning what (not) to forget,''
in \textit{Proc. ECCV}, 2018, pp.~139--154.

\bibitem{lopezpaz2017gradient}
D. Lopez-Paz and M. Ranzato,
``Gradient episodic memory for continual learning,''
in \textit{Proc. NeurIPS}, 2017, pp.~6467--6476.

\bibitem{chaudhry2019efficient}
A. Chaudhry, M. Ranzato, M. Rohrbach, and M. Elhoseiny,
``Efficient lifelong learning with A-GEM,''
in \textit{Proc. ICLR}, 2019.

\bibitem{rolnick2019experience}
D. Rolnick, A. Ahuja, J. Schwarz, T. Lillicrap, and G. Wayne,
``Experience replay for continual learning,''
in \textit{Proc. NeurIPS}, 2019.

\bibitem{mallya2018packnet}
A. Mallya and S. Lazebnik,
``PackNet: Adding multiple tasks to a single network by iterative pruning,''
in \textit{Proc. CVPR}, 2018, pp.~7765--7773.

\bibitem{rusu2016progressive}
A. A. Rusu et al.,
``Progressive neural networks,''
\textit{arXiv preprint arXiv:1606.04671}, 2016.

\bibitem{serra2018overcoming}
J. Serra, D. Suris, M. Miron, and A. Karatzoglou,
``Overcoming catastrophic forgetting with hard attention to the task,''
in \textit{Proc. ICML}, 2018, pp.~4548--4557.

\bibitem{amari1998natural}
S. Amari,
``Natural gradient works efficiently in learning,''
\textit{Neural Computation}, vol.~10, no.~2, pp.~251--276, 1998.

\bibitem{martens2015optimizing}
J. Martens and R. Grosse,
``Optimizing neural networks with Kronecker-factored approximate curvature,''
in \textit{Proc. ICML}, 2015, pp.~2408--2417.

\bibitem{krizhevsky2009learning}
A. Krizhevsky,
``Learning multiple layers of features from tiny images,''
Technical Report, University of Toronto, 2009.

\bibitem{mollahosseini2019affectnet}
A. Mollahosseini, B. Hasani, and M. H. Mahoor,
``AffectNet: A database for facial expression, valence, and arousal computing in the wild,''
\textit{IEEE Transactions on Affective Computing}, vol.~10, no.~1, pp.~18--31, 2019.

\bibitem{pearl2009causality}
J. Pearl,
\textit{Causality: Models, Reasoning and Inference}, 2nd~ed.
Cambridge University Press, 2009.

\end{thebibliography}

\end{document}
